% GENERATED FILE. Edit canonical lessons, then run npm run generate:books.
% canonical-source-hash: fnv1a64:82aa0ecdd79755bf
% canonical-lessons: KA-C442-heddari, KA-C442-tol, KA-C442-melsetuve, KA-C442-ghati, KA-C442-marga

\chapter{On the Highway}
\label{ch:ka-on-the-highway}
\hlchaptermodality{442}

\begin{chapteropening}
\textbf{By the end of this chapter:} \emph{I can say a highway, a toll, a flyover, a mountain pass and a route.}

Complete the last lesson of chapter 442: I can say a highway, a toll, a flyover, a mountain pass and a route.
\end{chapteropening}

\section[\texorpdfstring{\kn{ಹೆದ್ದಾರಿ} (heddāri)}{ಹೆದ್ದಾರಿ (heddāri)}]{\kn{ಹೆದ್ದಾರಿ} (heddāri) — a highway}

\label{lesson:KA-C442-heddari}



\begin{quote}
Before the new one: say the Kannada for cash, then the Kannada for a card.
\end{quote}

\subsection*{You'll want to know: \kn{ಹೆದ್ದಾರಿ}}
\textbf{\kn{ಹೆದ್ದಾರಿ}} — \emph{heddāri} — \textquotedblleft{}a highway\textquotedblright{}.

\begin{quote}
\kn{ಹೆದ್ದಾರಿಯಲ್ಲಿ} \kn{ವಾಹನಗಳು} \kn{ವೇಗವಾಗಿ} \kn{ಓಡುತ್ತವೆ}.
\end{quote}

\begin{quote}
\emph{heddāriyalli vāhanagaḷu vēgavāgi ōḍuttave.}
\end{quote}

\begin{quote}
Vehicles go fast on the highway.
\end{quote}

\kn{ಹೆ}- is the old big, as in \kn{ಹೆಬ್ಬೆರಳು}, thumb: the big road.

\begin{cousinweb}
\textbf{\kn{ಹೆದ್ದಾರಿ}} (\emph{heddāri}) is a native Kannada compound of \emph{per}, \emph{hē-}, \textquotedblleft{}big, great\textquotedblright{} (DEDR 4411), and \kn{ದಾರಿ}, \textquotedblleft{}a way, a road\textquotedblright{} (DEDR 3170).
\end{cousinweb}

\begin{grammarlens}[title={What you've built}]
One more word for getting around.
\end{grammarlens}

\subsection*{Your turn}
\begin{itemize}
\raggedright
  \item \emph{Say it:} \emph{heddāri}
  \item \emph{Say it:} \emph{heddāri}, once more
  \item \emph{Say it:} \emph{heddāriyalli vāhanagaḷu vēgavāgi ōḍuttave.} — Vehicles go fast on the highway.
  \item \emph{Recall:} say the Kannada for cash, then the Kannada for a card, then say \emph{heddāri} again
\end{itemize}

\begin{tcolorbox}[breakable,colback=teal!4,colframe=teal!35!black,title={Before you move on}]
What does \kn{ಹೆದ್ದಾರಿ} mean? (\textquotedblleft{}A highway\textquotedblright{}.) Say it once more.
\end{tcolorbox}
\section[\texorpdfstring{\kn{ಟೋಲ್} (ṭōl)}{ಟೋಲ್ (ṭōl)}]{\kn{ಟೋಲ್} (ṭōl) — a toll}

\label{lesson:KA-C442-tol}



\begin{quote}
Before the new one: say the Kannada for a card, then the Kannada for a highway.
\end{quote}

\subsection*{You'll want to know: \kn{ಟೋಲ್}}
\textbf{\kn{ಟೋಲ್}} — \emph{ṭōl} — \textquotedblleft{}a toll\textquotedblright{}.

\begin{quote}
\kn{ಹೆದ್ದಾರಿಯಲ್ಲಿ} \kn{ನಾವು} \kn{ಎರಡು} \kn{ಕಡೆ} \kn{ಟೋಲ್} \kn{ಕಟ್ಟಿದೆವು}.
\end{quote}

\begin{quote}
\emph{heddāriyalli nāvu eraḍu kaḍe ṭōl kaṭṭidevu.}
\end{quote}

\begin{quote}
We paid a toll at two places on the highway.
\end{quote}

\kn{ಟೋಲ್} \kn{ಕಟ್ಟು}, to pay a toll; \kn{ಟೋಲ್} \kn{ಗೇಟ್}, a toll gate.

\begin{cousinweb}
\textbf{\kn{ಟೋಲ್}} (\emph{ṭōl}) is a loan from English \emph{toll}, a charge for using a road or a bridge.
\end{cousinweb}

\begin{grammarlens}[title={What you've built}]
One more word for getting around.
\end{grammarlens}

\subsection*{Your turn}
\begin{itemize}
\raggedright
  \item \emph{Say it:} \emph{ṭōl}
  \item \emph{Say it:} \emph{ṭōl}, once more
  \item \emph{Say it:} \emph{heddāriyalli nāvu eraḍu kaḍe ṭōl kaṭṭidevu.} — We paid a toll at two places on the highway.
  \item \emph{Recall:} say the Kannada for a card, then the Kannada for a highway, then say \emph{ṭōl} again
\end{itemize}

\begin{tcolorbox}[breakable,colback=teal!4,colframe=teal!35!black,title={Before you move on}]
What does \kn{ಟೋಲ್} mean? (\textquotedblleft{}A toll\textquotedblright{}.) Say it once more.
\end{tcolorbox}
\section[\texorpdfstring{\kn{ಮೇಲ್ಸೇತುವೆ} (mēlsētuve)}{ಮೇಲ್ಸೇತುವೆ (mēlsētuve)}]{\kn{ಮೇಲ್ಸೇತುವೆ} (mēlsētuve) — a flyover}

\label{lesson:KA-C442-melsetuve}



\begin{quote}
Before the new one: say the Kannada for a highway, then the Kannada for a toll.
\end{quote}

\subsection*{You'll want to know: \kn{ಮೇಲ್ಸೇತುವೆ}}
\textbf{\kn{ಮೇಲ್ಸೇತುವೆ}} — \emph{mēlsētuve} — \textquotedblleft{}a flyover\textquotedblright{}.

\begin{quote}
\kn{ಹೊಸ} \kn{ಮೇಲ್ಸೇತುವೆಯ} \kn{ಮೇಲೆ} \kn{ಕಾರುಗಳು} \kn{ಓಡುತ್ತವೆ}.
\end{quote}

\begin{quote}
\emph{hosa mēlsētuveya mēle kārugaḷu ōḍuttave.}
\end{quote}

\begin{quote}
Cars run along the new flyover.
\end{quote}

\kn{ಮೇಲ್}, above, plus \kn{ಸೇತುವೆ}, a bridge: a bridge over a road.

\begin{cousinweb}
\textbf{\kn{ಮೇಲ್ಸೇತುವೆ}} (\emph{mēlsētuve}) is a Kannada compound of native \emph{mēl}, \textquotedblleft{}upper, above\textquotedblright{}, and \kn{ಸೇತುವೆ}, \textquotedblleft{}a bridge\textquotedblright{}, which is built on Sanskrit \emph{sētu}, \textquotedblleft{}a bond, a dam, a bridge\textquotedblright{} (Monier-Williams).
\end{cousinweb}

\begin{grammarlens}[title={What you've built}]
One more word for getting around.
\end{grammarlens}

\subsection*{Your turn}
\begin{itemize}
\raggedright
  \item \emph{Say it:} \emph{mēlsētuve}
  \item \emph{Say it:} \emph{mēlsētuve}, once more
  \item \emph{Say it:} \emph{hosa mēlsētuveya mēle kārugaḷu ōḍuttave.} — Cars run along the new flyover.
  \item \emph{Recall:} say the Kannada for a highway, then the Kannada for a toll, then say \emph{mēlsētuve} again
\end{itemize}

\begin{tcolorbox}[breakable,colback=teal!4,colframe=teal!35!black,title={Before you move on}]
What does \kn{ಮೇಲ್ಸೇತುವೆ} mean? (\textquotedblleft{}A flyover\textquotedblright{}.) Say it once more.
\end{tcolorbox}
\section[\texorpdfstring{\kn{ಘಾಟಿ} (ghāṭi)}{ಘಾಟಿ (ghāṭi)}]{\kn{ಘಾಟಿ} (ghāṭi) — a mountain pass; a ghat road}

\label{lesson:KA-C442-ghati}



\begin{quote}
Before the new one: say the Kannada for a toll, then the Kannada for a flyover.
\end{quote}

\subsection*{You'll want to know: \kn{ಘಾಟಿ}}
\textbf{\kn{ಘಾಟಿ}} — \emph{ghāṭi} — \textquotedblleft{}a mountain pass; a ghat road\textquotedblright{}.

\begin{quote}
\kn{ಘಾಟಿಯಲ್ಲಿ} \kn{ತುಂಬಾ} \kn{ತಿರುವುಗಳಿವೆ}.
\end{quote}

\begin{quote}
\emph{ghāṭiyalli tumbā tiruvugaḷive.}
\end{quote}

\begin{quote}
There are many bends on the ghat road.
\end{quote}

\kn{ಘಾಟಿ} \kn{ರಸ್ತೆ}, a ghat road, winds up through the hills.

\begin{cousinweb}
\textbf{\kn{ಘಾಟಿ}} (\emph{ghāṭi}) is usually traced to Sanskrit \emph{ghaṭṭa}, \textquotedblleft{}a landing place, a ferry\textquotedblright{}, whose basic sense, Turner notes, is a passage through mountains or over a river (CDIAL 4414).
\end{cousinweb}

\begin{grammarlens}[title={What you've built}]
One more word for getting around.
\end{grammarlens}

\subsection*{Your turn}
\begin{itemize}
\raggedright
  \item \emph{Say it:} \emph{ghāṭi}
  \item \emph{Say it:} \emph{ghāṭi}, once more
  \item \emph{Say it:} \emph{ghāṭiyalli tumbā tiruvugaḷive.} — There are many bends on the ghat road.
  \item \emph{Recall:} say the Kannada for a toll, then the Kannada for a flyover, then say \emph{ghāṭi} again
\end{itemize}

\begin{tcolorbox}[breakable,colback=teal!4,colframe=teal!35!black,title={Before you move on}]
What does \kn{ಘಾಟಿ} mean? (\textquotedblleft{}A mountain pass; a ghat road\textquotedblright{}.) Say it once more.
\end{tcolorbox}
\section[\texorpdfstring{\kn{ಮಾರ್ಗ} (mārga)}{ಮಾರ್ಗ (mārga)}]{\kn{ಮಾರ್ಗ} (mārga) — a route; a way}

\label{lesson:KA-C442-marga}



\begin{quote}
Before the new one: say the Kannada for a flyover, then the Kannada for a mountain pass.
\end{quote}

\subsection*{You'll want to know: \kn{ಮಾರ್ಗ}}
\textbf{\kn{ಮಾರ್ಗ}} — \emph{mārga} — \textquotedblleft{}a route; a way\textquotedblright{}.

\begin{quote}
\kn{ಈ} \kn{ಬಸ್ಸಿನ} \kn{ಮಾರ್ಗ} \kn{ಬದಲಾಯಿತು}.
\end{quote}

\begin{quote}
\emph{ī bassina mārga badalāyitu.}
\end{quote}

\begin{quote}
The route of this bus has changed.
\end{quote}

\kn{ಮಾರ್ಗ} is the formal word; \kn{ದಾರಿ} is the everyday way or path.

\begin{cousinweb}
\textbf{\kn{ಮಾರ್ಗ}} (\emph{mārga}) is borrowed from Sanskrit \emph{mārga}, \textquotedblleft{}a track, a road\textquotedblright{} (CDIAL 10071), first the trail of a wild animal.
\end{cousinweb}

\begin{grammarlens}[title={What you've built}]
One more word for getting around.
\end{grammarlens}

\subsection*{Your turn}
\begin{itemize}
\raggedright
  \item \emph{Say it:} \emph{mārga}
  \item \emph{Say it:} \emph{mārga}, once more
  \item \emph{Say it:} \emph{ī bassina mārga badalāyitu.} — The route of this bus has changed.
  \item \emph{Recall:} say the Kannada for a flyover, then the Kannada for a mountain pass, then say \emph{mārga} again
\end{itemize}

\begin{tcolorbox}[breakable,colback=teal!4,colframe=teal!35!black,title={Before you move on}]
What does \kn{ಮಾರ್ಗ} mean? (\textquotedblleft{}A route; a way\textquotedblright{}.) Say it once more.
\end{tcolorbox}
